\section{Limitations}
Text coverage is incomplete and uneven across datasets, which constrains every member-derived representation. The linked Problem Bodies subset is not the full KT benchmark, so results may differ when complete item text is available.

Several Chapter 5 values are thesis-text anchors because the raw artifacts do not always expose the exact validation-selection details needed to reconstruct every thesis table. We therefore use those values as sourced consolidation points rather than newly estimated statistics.

The functional graph analyses are observational. Transition, co-occurrence, difficulty-similarity, and DKT-derived graphs are proxies for functional structure, not definitive prerequisite graphs. The datasets are math-heavy ASSISTments benchmarks, so generalization to other domains and instructional contexts remains open.

Finally, this paper studies fixed KC labels. It does not evaluate methods that learn new category boundaries. Chapter 6-style category redefinition is a separate research question and should not be treated as evidence for the claims in this short paper.
